\documentclass[11pt,a4paper]{article}
\usepackage[T1]{fontenc}
\usepackage{lmodern,amsmath,amssymb,amsthm,mathtools}
\usepackage[margin=25mm]{geometry}
\usepackage{xcolor}
\usepackage[colorlinks=true,linkcolor=blue!45!black,citecolor=blue!45!black,urlcolor=blue!45!black]{hyperref}
\usepackage{microtype}
\newtheorem{theorem}{Theorem}[section]
\newtheorem{proposition}[theorem]{Proposition}
\newtheorem{lemma}[theorem]{Lemma}
\newtheorem{corollary}[theorem]{Corollary}
\theoremstyle{definition}
\newtheorem{definition}[theorem]{Definition}
\newtheorem{remark}[theorem]{Remark}
\newcommand{\OKS}{\mathcal O_{K,S}}
\newcommand{\A}{\mathbb A}
\newcommand{\PP}{\mathbb P}
\setlength{\emergencystretch}{2em}
\hypersetup{pdftitle={Degree and variable lower bounds for Diophantine definitions},pdfsubject={Unconditional obstructions over number fields and rings of S-integers}}
\title{Degree and Variable Lower Bounds\\for Diophantine Definitions}
\author{}
\date{9 November 2026}
\begin{document}
\maketitle
\begin{abstract}
Let $K$ be a number field and let $S$ be a finite set of its finite
primes. A single polynomial whose $K$-rational input projection is
an infinite subset of the ring of $S$-integers must have total
degree at least four and at least three total coordinates. The
coordinate bound also holds for finite systems of arbitrary degree,
even with nonzero conditions and finite alternatives. Systems of
quadratic equations require at least five total coordinates; systems
of equations of degree at most three require at least four. More
precisely, a system in three coordinates with an
infinite integral projection must have a common irreducible factor
of degree at least four. For a single polynomial, the input and every essential witness coordinate must
each have individual degree at least two. The proof combines an
elementary argument about prime divisors of polynomial values,
Siegel's theorem for integral points on curves, and Koll\'ar's
unirationality theorem for cubic hypersurfaces. Projection from a
smooth point reduces a surface intersection of two quadrics to a
cubic surface, including the necessary reducible and nonreduced
cases. Every nonnegative quartic requires at least five coordinates,
and a nonnegative polynomial in three coordinates requires
degree at least eight. These are necessary coordinate bounds;
attainment of the unrestricted three-coordinate bound is not established.
\end{abstract}

\section{Definitions and statements}

For a nonzero prime ideal $\mathfrak p$ of the ring of integers
$\mathcal O_K$, write $v_{\mathfrak p}$ for its additive valuation,
normalized to have value group $\mathbb Z$, and put
$v_{\mathfrak p}(0)=+\infty$. For a finite set $S$ of nonzero
prime ideals, define
\[
 \OKS=\{a\in K: v_{\mathfrak p}(a)\ge0
                       \text{ for every }\mathfrak p\notin S\}.
\]
In particular, $\mathcal O_{K,\varnothing}=\mathcal O_K$.

\begin{definition}
A single-polynomial presentation over $K$ is an equation
\[
                P(T,Y_1,\ldots,Y_m)=0,
       \qquad P\in K[T,Y_1,\ldots,Y_m].
\]
Its \emph{input projection} is the set
\[
 A_P=\{t\in K:\exists \mathbf y\in K^m, P(t,\mathbf y)=0\}.
\]
The distinguished input $T$ is counted among the coordinates; thus
there are $m+1$ total coordinates and $m$ existential witnesses.
Degree means total degree, including $T$. A witness coordinate is
\emph{essential} if the polynomial depends on it.
\end{definition}

A finite system uses the same coordinates for all its equations.
We also consider finite unions of presentations consisting of equations
and nonzero conditions. Geometrically, each such presentation is a
Zariski open subset of a common zero set. All coordinates are evaluated
in $K$, and no ordering conditions or universal quantifiers are allowed.

\begin{samepage}
\begin{theorem}\label{lb:main}
Let $K$ be a number field and $S$ a finite set of finite primes.
\begin{enumerate}
\item If a single-polynomial presentation has an infinite input
projection contained in $\OKS$, then it has at least three total
coordinates and total degree at least four.
\item The same lower bound of three total coordinates holds for
finite systems of arbitrary degree, with nonzero conditions and finite
alternatives permitted.
\item In the single-polynomial case, the individual degree in $T$
and in each essential witness coordinate is at least two.
\item A finite system whose equations have degree at most three
and whose input projection is infinite and contained in $\OKS$
requires at least four total coordinates, even with nonzero
conditions and finite alternatives.
\item A finite system of quadratic equations with an infinite
input projection contained in $\OKS$ requires at least five total
coordinates. Nonzero conditions and finite alternatives are permitted.
\end{enumerate}
\end{theorem}
\end{samepage}

For $K=\mathbb Q$ and $S=\varnothing$, this applies in particular
to every polynomial definition of $\mathbb Z$ over $\mathbb Q$.
The degree restriction in part~(1) concerns one polynomial. A system
of quadratic equations is a different presentation, and is not
excluded by that restriction. In particular, part~(5) applies to
quartics written as positive sums of squares of quadratic equations.
Section~\ref{lb:positivity-section} extends the five-coordinate bound
to every nonnegative quartic over a field with a real embedding.
The bound is not asserted here for every quartic.

Two established geometric results enter the proof. Siegel's theorem
gives finiteness of integral points on affine curves of positive
genus, including $S$-integral points over a number field
\cite[Chapter~8]{Lang}. Koll\'ar's theorem gives unirationality of a
noncone cubic hypersurface of dimension at least two over a perfect
field as soon as it has a rational point \cite[Theorem~2]{Kollar}.
All number fields are perfect. The other ingredients are proved below.

\section{Rational functions and primes outside a finite set}

\begin{lemma}\label{lb:prime-divisors}
Let $g\in\mathcal O_K[Z]$ be nonconstant. For every finite set
$\Sigma$ of nonzero prime ideals of $\mathcal O_K$, there exist
$a\in\mathcal O_K$ and $\mathfrak p\notin\Sigma$ such that
\[
                  g(a)\ne0,
          \qquad v_{\mathfrak p}(g(a))>0.
\]
\end{lemma}
\begin{proof}
After replacing $g(Z)$ by $g(Z+c)$ for a suitable $c\in\mathcal O_K$,
we may suppose that $b=g(0)\ne0$. Choose a nonzero rational integer
$M$ belonging to every ideal in $\Sigma$; take $M=1$ if $\Sigma$ is
empty. Write $g(Z)=b+c_1Z+\cdots+c_dZ^d$, where $d\ge1$ and $c_d\ne0$.
Then
\[
 h(N)=\frac{g(bMN)}b
      =1+c_1MN+c_2bM^2N^2+\cdots+c_db^{d-1}M^dN^d
      \in\mathcal O_K[N].
\]
For every rational integer $n$ and every $\mathfrak q\in\Sigma$,
we have $h(n)\equiv1\pmod{\mathfrak q}$. The norm polynomial
\[
                    \prod_{\sigma:K\hookrightarrow\mathbb C}h^\sigma(N)
\]
lies in $\mathbb Z[N]$ and has degree $d[K:\mathbb Q]$: its
leading coefficient is the nonzero norm of $c_db^{d-1}M^d$.
Its absolute value is therefore greater than one for all sufficiently
large positive integers $n$. Such an $h(n)$ is nonzero and is not a
unit of $\mathcal O_K$. A prime ideal dividing it lies outside
$\Sigma$ and also divides the nonzero value $g(bMn)=b h(n)$.
Undoing the initial translation gives the required $a$.
\end{proof}

\begin{lemma}\label{lb:rational-function}
Let $r\in K(Z)$ be nonconstant, and let $E\subset K$ be finite.
There is an $a\in K\setminus E$ at which $r$ is defined and
$r(a)\notin\OKS$.
\end{lemma}
\begin{proof}
First suppose that $r$ is a polynomial. Choose a prime ideal
$\mathfrak p\notin S$ at which all its coefficients are integral
and its leading coefficient is a unit. Only finitely many primes
are excluded by these conditions. Choose $0\ne\pi\in\mathfrak p$.
For $a=\pi^{-n}$ and sufficiently large $n$, the leading term of
$r(a)$ has strictly smaller $\mathfrak p$-adic valuation than all
other terms. Its valuation is negative, so
$v_{\mathfrak p}(r(a))<0$. The numbers $\pi^{-n}$ are distinct,
and hence all but finitely many avoid $E$.

Otherwise write $r=f/g$ with relatively prime polynomials
$f,g\in\mathcal O_K[Z]$ and with $g$ nonconstant. The polynomial
Euclidean algorithm over $K$ and clearing denominators give
$A,B\in\mathcal O_K[Z]$ and $0\ne c\in\mathcal O_K$ satisfying
\[
                              Af+Bg=c.
\]
Enlarge $S$ to a finite set $\Sigma$ containing the prime ideals
dividing $c$ or any nonzero value $g(e)$ with
$e\in E\cap\mathcal O_K$.
By Lemma~\ref{lb:prime-divisors}, there are $a\in\mathcal O_K$ and
$\mathfrak p\notin\Sigma$ with $g(a)\ne0$ and
$v_{\mathfrak p}(g(a))>0$. The displayed identity implies
$v_{\mathfrak p}(f(a))=0$. Thus $r(a)$ has negative valuation at
$\mathfrak p\notin S$. Moreover $a\notin E$: if $a=e$ were in
$E\cap\mathcal O_K$, the prime divisor $\mathfrak p$ of its
nonzero value $g(e)$ would have been included in $\Sigma$.
\end{proof}

An integral $K$-variety is called \emph{$K$-unirational} if it is
dominated by an affine space through a rational map defined over $K$.

\begin{proposition}\label{lb:unirational}
Let $X$ be a $K$-unirational integral variety, let $U\subseteq X$
be nonempty and open, and let $f:U\to\A^1_K$ be regular. If
\[
                              f(U(K))\subseteq\OKS,
\]
then $f$ is constant.
\end{proposition}
\begin{proof}
Pull back $f$ through a dominant rational parametrization. If $f$
is nonconstant, the result is a nonconstant rational function on
affine space. Shrink its permitted domain to a nonempty principal
open subset. In characteristic zero, one of its partial derivatives
is nonzero. Since $K$ is infinite, we can fix the other coordinates
in $K$ so that this derivative remains nonzero and none of the
polynomials defining the permitted open set vanishes identically
on the resulting line. The restriction is a nonconstant rational
function $r\in K(Z)$, and the parametrization is permitted at every
$K$-point of the line except a finite set. Lemma~\ref{lb:rational-function}
then supplies a value outside $\OKS$, a contradiction.
\end{proof}

\section{Integral-valued functions on curves}

\begin{lemma}\label{lb:positive-genus}
Let $C$ be a smooth projective geometrically integral curve over
$K$ of positive genus, and let $f\in K(C)$ be nonconstant. There
are only finitely many $P\in C(K)$ at which $f$ is defined and
$f(P)\in\OKS$.
\end{lemma}
\begin{proof}
The function $f$ defines a finite morphism $C\to\PP^1_K$.
Removing its poles gives an affine curve
$D=f^{-1}(\A^1_K)$ finite over $\A^1_K$. Choose generators
$z_1,\ldots,z_r$ of its coordinate ring over $K[f]$. Each $z_i$
satisfies a monic equation over $K[f]$.

There is a finite set $S'$ of finite primes, containing $S$, such
that all coefficients in these finitely many equations belong to
$\mathcal O_{K,S'}[f]$. If $P\in D(K)$ and $f(P)\in\OKS$,
then each $z_i(P)\in K$ is integral over $\mathcal O_{K,S'}$.
That ring is integrally closed, so all $z_i(P)$ belong to it.
Consequently these points are $S'$-integral in a fixed affine
embedding of $D$. Siegel's theorem gives their finiteness
\cite[Chapter~8]{Lang}. Its positive-genus hypothesis holds for
$D$, and at least one pole was removed because $f$ is nonconstant.
\end{proof}

\begin{proposition}\label{lb:curves}
Let $X$ be an algebraic set over $K$ of dimension at most one,
let $U\subseteq X$ be open, and let $f:U\to\A^1_K$ be regular.
If $f(U(K))\subseteq\OKS$, then $f(U(K))$ is finite.
\end{proposition}
\begin{proof}
Pass to the finitely many irreducible components, omitting components
disjoint from $U$. Zero-dimensional components contribute finitely
many points. If an irreducible curve is not geometrically integral,
every $K$-point lies in the intersection of distinct conjugate
geometric components. These intersections form a finite set, so this
case also contributes only finitely many points.

For a geometrically integral curve, take its smooth projective
normalization $C$. Away from finitely many points the normalization
is an isomorphism, and $f$ pulls back to a rational function on $C$.
A constant function has finite image. If the genus is positive,
Lemma~\ref{lb:positive-genus} applies. If the genus is zero and
$C(K)$ is empty, only the finitely many points omitted in passing
to the normalization can contribute. If the genus is zero and
$C(K)$ is nonempty, then $C\simeq\PP^1_K$; Proposition~\ref{lb:unirational}
excludes a nonconstant function with the asserted integral-value
property on the permitted open subset. This proves the result.
\end{proof}

\section{Hypersurfaces of degree at most three}

\begin{theorem}\label{lb:cubic}
Let $X\subseteq\PP^N_K$ be the zero set of a nonzero homogeneous
polynomial of degree at most three. For every open $U\subseteq X$
and regular function $f:U\to\A^1_K$, the containment
$f(U(K))\subseteq\OKS$ implies that $f(U(K))$ is finite.
The same conclusion holds for $X=\PP^N_K$.
\end{theorem}
\begin{proof}
Projective space is unirational, so that case follows from
Proposition~\ref{lb:unirational}. Factor the equation of a hypersurface
over $K$ and treat each irreducible factor separately. Multiplicities
do not change the rational point set.

\emph{Geometrically reducible components.}
An irreducible polynomial of degree two or three over a perfect
field, if not geometrically irreducible, splits geometrically into
conjugate linear factors. Indeed the geometric factors have a common
degree and form one Galois orbit; that degree divides the prime
number two or three. Every $K$-point belongs to all the conjugate
hyperplanes. Their common intersection is a $K$-linear subspace,
and Proposition~\ref{lb:unirational} applies to the restriction of $f$.
Linear components are treated by the same proposition.

\emph{Geometrically integral quadrics.}
A quadric with a smooth $K$-point is rational by projection from
that point. If it has no smooth $K$-point, all its $K$-points lie
in its singular locus, which is a linear subspace. Both cases
give the required finite image.

\emph{Cubics that are not cones.}
Curves and zero-dimensional cases are covered by
Proposition~\ref{lb:curves}. A geometrically integral cubic of
dimension at least two that is not a cone is $K$-unirational if
it has a $K$-point, by Koll\'ar's theorem \cite[Theorem~2]{Kollar}.
This theorem applies to singular cubics as well as smooth ones.
If there is no $K$-point the image is empty; otherwise use
Proposition~\ref{lb:unirational}.

\emph{The structure of a cubic cone.}
Write its cubic equation as $F$ on a $K$-vector space $V$ and put
\[
 W=\{w\in V:\partial_wF=0\text{ as a polynomial}\}.
\]
This is the kernel of a $K$-linear map. In characteristic zero,
$F$ is constant along the directions in $W$, so it descends to
a cubic form on $V/W$. The projective space $\PP(W)$ is the
maximal vertex. Choose a $K$-linear complement to $W$. It gives
a geometrically integral cubic base $B$ which is not a cone.
Indeed, a vertex direction of the base would lift to an additional
direction in the kernel defining $W$. That kernel commutes with
extension of the ground field, so the same assertion holds geometrically.
Its dimension is at least one, since a zero-dimensional cubic
hypersurface is not geometrically integral.

Outside $\PP(W)$, projection from the vertex gives a morphism
to $B$. Each fiber over a $K$-point of $B$ is an affine space of
dimension $\dim W$. On an open subset of $B$ this bundle is
trivial, so the cone is birational to $B\times\A^{\dim W}$.
These statements follow directly by writing a point of $V$ as
the sum of a vector in $W$ and a vector in the chosen complement.

If $B(K)$ is empty, all rational points of the cone lie in the
vertex. If $B(K)$ is finite, apply Proposition~\ref{lb:unirational}
on the permitted open part of each affine-space fiber, and also
on the vertex. This gives a finite union of finite images.
If $\dim B\ge2$ and $B(K)$ is nonempty, Koll\'ar's theorem
makes $B$, and then the cone, unirational.

\emph{A cone over a curve with infinitely many rational points.}
The remaining case has $\dim B=1$ and $B(K)$ infinite. Let $C$
be the smooth projective normalization of $B$. Discarding finitely
many base points, write the function $f$ on an open subset of
$C\times\A^{\dim W}$. If this rational function depends
nontrivially on the affine-space coordinates, some partial
derivative in those coordinates is nonzero. Expressing it as a
quotient of polynomials with coefficients in $K(C)$ shows that,
outside finitely many base points, the specialized function is
nonconstant and its permitted fiber is nonempty. Choose a
$K$-point of $C$ outside those exceptions. Its fiber is a
nonempty open subset of affine space, contradicting
Proposition~\ref{lb:unirational}. Therefore $f$ belongs to $K(C)$.

For every $K$-point of $C$ outside a finite set, the permitted
fiber is a nonempty open subset of affine space and has a $K$-point.
Thus the descended function is defined and $S$-integral at all
these base points. Proposition~\ref{lb:curves} implies that its
image there is finite. Since a nonconstant rational function on
a curve has finite fibers and $C(K)$ is infinite, the function
must be constant. It is then constant as an element of the
cone's function field, and hence wherever it is regular on that
irreducible cone. This includes the points omitted from the
chosen birational chart.
\end{proof}

\section{Consequences for polynomial presentations}

\begin{proof}[Proof of Theorem~\ref{lb:main}(1) and (2)]
First consider at most two total coordinates. For one system of
equations, its zero set is either all of the ambient affine space
or has dimension at most one. In the first case, the permitted
nonempty open subset is unirational, and its input function cannot
have an infinite $S$-integral image by Proposition~\ref{lb:unirational}.
In the second case, Proposition~\ref{lb:curves} gives a finite
input image. Empty permitted sets contribute nothing. Taking
finitely many alternatives preserves finiteness. This proves the
coordinate bound for systems and for single equations.

For a single nonconstant polynomial of total degree at most three,
homogenize it and regard the distinguished input as a regular
function on the usual affine chart of its projective zero set.
Theorem~\ref{lb:cubic} gives a finite input image, independently
of the number of coordinates. A nonzero constant polynomial has
empty projection. The zero polynomial has projection $K$, which
is not contained in $\OKS$ because $S$ is finite. These observations
prove the degree bound.
\end{proof}

\begin{proposition}\label{lb:individual-degrees}
Suppose that $P\in K[T,Y_1,\ldots,Y_m]$ has an infinite input
projection contained in $\OKS$. Then
\[
                \deg_T P\ge2,
    \qquad \deg_{Y_j}P\ge2
          \quad\text{for every essential witness }Y_j.
\]
\end{proposition}
\begin{proof}
Suppose first that $P$ is linear in an essential witness $Y_j$.
Write
\[
             P=A(T,\mathbf Y')Y_j+B(T,\mathbf Y'),\qquad A\ne0.
\]
Choose constants $\mathbf c\in K^{m-1}$ so that
$A(T,\mathbf c)$ is a nonzero polynomial. Such a choice exists
because $K$ is infinite. For every $t$ outside its finite root
set, setting $Y_j=-B(t,\mathbf c)/A(t,\mathbf c)$ gives a
rational witness. The projection is therefore cofinite in $K$.
It cannot be contained in $\OKS$, whose complement is infinite:
for a prime $\mathfrak p\notin S$ and $0\ne\pi\in\mathfrak p$,
all the distinct elements $\pi^{-n}$ lie outside $\OKS$.

Suppose next that $\deg_TP\le1$, and write
$P=A(\mathbf Y)T+B(\mathbf Y)$. If $A=0$, the input projection
is either empty or all of $K$. Otherwise on the nonempty open
set $A\ne0$ the rational function $-B/A$ takes only values in
$\OKS$. Proposition~\ref{lb:unirational} makes it constant.
At any rational point with $A=B=0$, every input would be admitted,
which is impossible. Thus the projection has at most one element,
again a contradiction. This proves Theorem~\ref{lb:main}(3).
\end{proof}

\begin{corollary}\label{lb:linear-systems}
No finite system of linear equations, even with nonzero conditions
and finite alternatives, has an infinite input projection contained
in $\OKS$.
\end{corollary}
\begin{proof}
Each nonempty common zero set of linear equations is an affine
space over $K$. Apply Proposition~\ref{lb:unirational} to its
permitted open subset and take the finite union of the resulting
finite images.
\end{proof}

\begin{proposition}\label{lb:three-coordinate-systems}
Let $F_1,\ldots,F_s\in K[T,Y_1,Y_2]$, and let $U$ be a
Zariski open subset of their common zero set. Suppose that
$T(U(K))$ is infinite and contained in $\OKS$. Then some
irreducible polynomial of degree at least four divides every
$F_i$. In particular every nonzero $F_i$ has degree at least four.
The nonzero conditions defining $U$ may have arbitrary degrees.
\end{proposition}
\begin{proof}
The case in which all $F_i$ vanish identically is impossible:
then $U$ is a nonempty open subset of affine three-space, and
Proposition~\ref{lb:unirational} would make its input function
constant. Discard any identically zero equations and put
$g=\gcd(F_1,\ldots,F_s)$, computed in the unique factorization
domain $K[T,Y_1,Y_2]$. Write $F_i=gh_i$. Then
\[
 V(F_1,\ldots,F_s)=V(g)\cup V(h_1,\ldots,h_s).
\]
The second zero set has dimension at most one. Indeed, an
irreducible component of dimension two would correspond to a
height-one prime of the polynomial ring. Such a prime is
generated by one irreducible polynomial, which would divide
all the $h_i$, contrary to $\gcd(h_1,\ldots,h_s)=1$.
A component of dimension three would make all the $h_i$ zero,
which is also impossible.

By Proposition~\ref{lb:curves}, the permitted open subset of
$V(h_1,\ldots,h_s)$ has finite input image. Factor $g$ into
irreducibles. There are finitely many resulting hypersurfaces,
so at least one factor $g_0$ satisfies
\[
                   T\bigl((U\cap V(g_0))(K)\bigr)
                       \quad\text{is infinite}.
\]
Every such input lies in $\OKS$. If $\deg g_0\le3$,
homogenizing $g_0$ and applying Theorem~\ref{lb:cubic} to
this permitted affine open subset gives a finite image,
a contradiction. Therefore $\deg g_0\ge4$. Since $g_0$
divides $g$, it divides every defining equation.
\end{proof}

\begin{corollary}\label{lb:low-degree-systems}
A finite system of equations of total degree at most three
with an infinite input projection contained in $\OKS$ uses
at least four total coordinates. The same conclusion holds
with arbitrary nonzero conditions and finitely many alternatives.
For quadratic systems, Section~\ref{lb:quadratic-section} gives the stronger bound of
five total coordinates.
\end{corollary}
\begin{proof}
If there are fewer than three coordinates, adjoining unused
coordinates gives a presentation in three coordinates with
the same input projection and the same equation degrees.
Proposition~\ref{lb:three-coordinate-systems} excludes an
infinite image for each alternative. A finite union of finite
images is finite. This proves Theorem~\ref{lb:main}(4).
\end{proof}

\section{Quadratic systems in four coordinates}\label{lb:quadratic-section}

The extra coordinate obstruction for quadratic systems follows from
an explicit projection. The multiplicity issue matters: the common
zero set of two quadrics need not be reduced.

\begin{lemma}\label{lb:small-degree-surface}
Let $Y\subseteq\PP^4_K$ be a geometrically integral projective
surface of degree at most three. For every open $V\subseteq Y$ and regular
function $f:V\to\A^1_K$, if $f(V(K))\subseteq\OKS$, then
$f(V(K))$ is finite.
\end{lemma}
\begin{proof}
The elementary degree inequality
\[
 \deg Y\ge 1+\operatorname{codim}(Y,\langle Y\rangle)
\]
for a projective integral variety is recalled in
\cite[Introduction, inequality $(*)$]{EGHP}. The linear span
$\langle Y\rangle$ is defined over $K$. If the span has dimension
at most three, $Y$ is a plane or a hypersurface of degree at most
three in that span, so Theorem~\ref{lb:cubic} applies.

The remaining case has $\deg Y=3$ and $\langle Y\rangle=\PP^4$.
If $V(K)$ is not dense, its closure has dimension at most one,
and Proposition~\ref{lb:curves} applies. Otherwise choose a smooth
point $P\in V(K)$. Projection from $P$ has a surface as its image:
if the image had dimension one, $Y$ would be a cone with vertex
$P$; every generating line would then lie in the projective
tangent plane at the smooth point $P$, contradicting the span.
Write $Y'\subseteq\PP^3$ for the image surface, and $e$ for the
generic degree of the projection. The image is nondegenerate,
because a hyperplane containing it would pull back to a
hyperplane containing $Y$. Intersecting with a general line in
$\PP^3$, whose inverse image is a plane through $P$, gives
\[
                         3=1+e\deg Y'.
\]
Here the contribution at $P$ is one because $P$ is smooth and
the plane is general. This is the point-projection degree
formula \cite[Theorem~5.11]{Mumford}, with multiplicity one.
A nondegenerate surface in $\PP^3$ has
degree at least two; thus $e=1$ and $Y'$ is a quadric.
On isomorphic nonempty open subsets the rational function $f$
therefore corresponds to a regular function on an open subset
of that quadric. Theorem~\ref{lb:cubic} gives finite image there.
The omitted subset of $Y$ has dimension at most one, and
Proposition~\ref{lb:curves} handles its image.
\end{proof}

\begin{proposition}\label{lb:two-quadrics}
Let $Q_1,Q_2\in K[X_0,\ldots,X_4]$ be relatively prime
homogeneous quadratic forms. Let $Y$ be an irreducible surface
component of their projective common zero set, endowed with its
reduced structure. If $V\subseteq Y$ is open and $f:V\to\A^1_K$
is regular with $f(V(K))\subseteq\OKS$, then $f(V(K))$ is finite.
\end{proposition}
\begin{proof}
If $Y$ is not geometrically integral, all its $K$-points lie in
the intersections of distinct conjugate geometric components.
Their union descends to a proper closed subset of $Y$ of dimension
at most one. Proposition~\ref{lb:curves} proves the assertion in
this case. Hence assume that $Y$ is geometrically integral.

If $V(K)$ is not Zariski dense in $Y$, its Zariski closure is
again of dimension at most one, and the same proposition applies.
Suppose therefore that $V(K)$ is dense. B\'ezout's theorem for the
complete intersection $Z=V(Q_1,Q_2)\subset\PP^4_K$ gives
\[
                    \sum_i m_i\deg Y_i=4,
\]
where $Y_i$ are the geometric surface components and $m_i$ their
multiplicities at their generic points. If the multiplicity of
$Y$ is at least two, then $\deg Y\le2$, and
Lemma~\ref{lb:small-degree-surface} applies.

It remains to treat multiplicity one. The complete intersection
$Z$ is then reduced at the generic point of $Y$. Over the perfect
field $K$, there is a nonempty open subset of $Y$ on which $Z$
is smooth of dimension two and no other component occurs. Density
supplies a point $P\in V(K)$ in that open subset. After a
$K$-linear projective change of coordinates, put
$P=(1:0:0:0:0)$. Because $Q_i(P)=0$, write
\[
 Q_1=X_0L_1(\mathbf X)+q_1(\mathbf X),\qquad
 Q_2=X_0L_2(\mathbf X)+q_2(\mathbf X),
 \qquad \mathbf X=(X_1,\ldots,X_4),
\]
where the $L_i$ are linear and the $q_i$ quadratic. Smoothness
of $Z$ at $P$ says that $L_1,L_2$ are linearly independent.

If both $L_i$ vanish identically on $Y$, then $Y$ is the plane
$V(L_1,L_2)$, and Lemma~\ref{lb:small-degree-surface} applies.
Otherwise, after interchanging the two equations, assume that
$L_1$ does not vanish identically on $Y$. Projection from $P$
then has the rational inverse
\[
 [X_1:\cdots:X_4]\longmapsto
 [-q_1(\mathbf X):L_1(\mathbf X)X_1:\cdots:
                               L_1(\mathbf X)X_4]
\]
on the open set $L_1\ne0$. The image lies in the cubic
\[
              H=L_1q_2-L_2q_1=0\quad\text{in }\PP^3_K.
\]
This polynomial is nonzero. Indeed, if $H=0$, relative
primality of the independent linear forms would give
$q_1=L_1M$ and $q_2=L_2M$ for one linear form $M$.
Both quadrics would then be divisible by $X_0+M$, contrary
to their relative primality.

The image closure of $Y$ is thus a surface component $Y'$ of
the cubic $V(H)$. It is defined by an irreducible factor of
$H$, of degree at most three. The displayed inverse gives an
isomorphism between nonempty open subsets of $Y$ and $Y'$.
Shrink the first open subset so that it lies in $V$. On the
corresponding open subset of $Y'$, the function $f$ is a
\emph{rational} function, regular after this shrinking, and
all its $K$-rational values lie in $\OKS$.
Theorem~\ref{lb:cubic} makes its image finite. The omitted
closed subset of the surface $Y$ has dimension at most one;
Proposition~\ref{lb:curves} makes the remaining image finite.
\end{proof}

\begin{theorem}\label{lb:quadratic-five}
Let $F_1,\ldots,F_s\in K[T,Y_1,Y_2,Y_3]$ have total degree
at most two, and let $U$ be open in their common zero set.
If $T(U(K))\subseteq\OKS$, then $T(U(K))$ is finite.
Consequently every finite quadratic system with infinite
$S$-integral input projection needs at least five total coordinates.
The assertion is unchanged by nonzero conditions or finitely many
alternatives.
\end{theorem}
\begin{proof}
Identically zero equations may be discarded. If all equations
are zero, use Proposition~\ref{lb:unirational}; a nonzero constant
equation gives the empty set. A nonconstant linear equation
allows a $K$-affine parametrization of its hyperplane by three
coordinates. If the input is constant on that hyperplane, the
claim is immediate. Otherwise choose the parametrization with
the input itself as one of those coordinates. The remaining
equations still have degree at most two, so
Corollary~\ref{lb:low-degree-systems} applies. We may therefore
assume that all remaining equations have degree exactly two.

If $g=\gcd(F_1,\ldots,F_s)$ is nonconstant, write $F_i=gh_i$.
Their common zero set is
\[
       V(g)\ \cup\ V(h_1,\ldots,h_s).
\]
The first part has finite permitted input image by
Theorem~\ref{lb:cubic}. The $h_i$ have degree at most one,
so the second part has finite input image by
Corollary~\ref{lb:linear-systems}.

Suppose now that the gcd is one. Choose one nonzero equation
$F$. A $K$-linear combination $G$ of the equations can be
chosen relatively prime to $F$: for each of the finitely many
irreducible factors of $F$, divisibility of the combination
by that factor is a proper linear condition on its coefficients.
A finite union of proper linear subspaces cannot cover the
coefficient space over the infinite field $K$. The case
$\deg G\le1$ is already settled, so take $\deg G=2$.

The homogenizations $Q_1,Q_2$ of $F,G$ in $\PP^4_K$ remain
relatively prime. A common factor other than the homogenizing
coordinate would dehomogenize to a common factor of $F,G$;
neither homogenization is divisible by that coordinate.
Their intersection is therefore a complete intersection of
pure dimension two. Every surface component of the projective
closure of $V(F_1,\ldots,F_s)$ is a component of this
intersection. Proposition~\ref{lb:two-quadrics} makes its
permitted input image finite. Components of dimension at most
one are covered by Proposition~\ref{lb:curves}. There are
only finitely many components, proving the assertion.

Adjoining unused coordinates treats systems in fewer than four
coordinates, and a finite union of finite images is finite.
This proves Theorem~\ref{lb:main}(5).
\end{proof}

\begin{corollary}\label{lb:quartic-sos}
Let $K$ have a real embedding $\sigma$, let $c_1,\ldots,c_s\in K$
satisfy $\sigma(c_i)>0$, and let $F_i\in K[T,\mathbf Y]$ have
degree at most two. If
\[
                  P=\sum_{i=1}^s c_iF_i^2
\]
has infinite input projection contained in $\OKS$, then $P$
uses at least five total coordinates.
\end{corollary}
\begin{proof}
For a $K$-rational tuple, $P=0$ holds if and only if every
$F_i=0$, as follows by applying $\sigma$. Apply
Theorem~\ref{lb:quadratic-five}.
\end{proof}

\section{Three-coordinate presentations and nonnegative polynomials}\label{lb:positivity-section}

\begin{proposition}\label{lb:three-surface}
Suppose that $P\in K[T,Y_1,Y_2]$ has infinite input projection
contained in $\OKS$. Then $P$ has an absolutely irreducible factor
$G$ of degree at least four such that $V(G)(K)$ is Zariski dense
and its input projection is infinite. Every proper closed subset
of $V(G)$ has finite input projection on its $K$-points.
In particular, if $\deg P=4$, then $P$ is absolutely irreducible
up to a nonzero constant factor and its smooth $K$-points have
infinite input projection.
\end{proposition}
\begin{proof}
Factor $P$ over $K$. Some irreducible factor $G$ has infinite
input projection, since a finite union of finite sets is finite.
If $G$ is not absolutely irreducible, its rational points lie
in the intersections of conjugate geometric components, a
closed subset of dimension at most one. Proposition~\ref{lb:curves}
would give a finite image. Hence $G$ is absolutely irreducible.
The same proposition shows that no proper closed subset of its
surface can contain all rational points contributing infinitely
many inputs. Thus the rational points are dense. The degree
bound follows from Theorem~\ref{lb:cubic}. If $\deg P=4$,
the factor $G$ already has the full degree and occurs once.
Its singular locus is a proper closed subset of dimension at
most one, hence contributes only finitely many inputs.
\end{proof}

\begin{theorem}\label{lb:nonnegative-eight}
Let $K$ have a real embedding $\sigma$. Suppose that
$P\in K[T,Y_1,Y_2]$ satisfies
\[
             P^\sigma(t,y_1,y_2)\ge0
                  \quad\text{for all }(t,y_1,y_2)\in\mathbb R^3.
\]
If its $K$-rational input projection is infinite and contained
in $\OKS$, then $P$ has a repeated absolutely irreducible factor
of degree at least four. Consequently $\deg P\ge8$.
\end{theorem}
\begin{proof}
Choose the factor $G$ of Proposition~\ref{lb:three-surface}.
Every real zero of a nonnegative differentiable function is a
local minimum, so all first partial derivatives of $P^\sigma$
vanish at every point of $V(G)(K)$ after applying $\sigma$.
The embedding is injective. Therefore all first partial
derivatives of $P$ vanish on the Zariski-dense set $V(G)(K)$,
and $G$ divides each of them.

Write $P=G^eH$ with $G\nmid H$. If $e=1$, reduction of
$\partial P/\partial X_j$ modulo $G$ gives
\[
               H\frac{\partial G}{\partial X_j}=0
                    \quad\text{in }K[T,Y_1,Y_2]/(G)
\]
for each coordinate $X_j$. This quotient is an integral domain
and $H$ is nonzero in it, so $G$ divides every partial derivative
of itself. Their degrees are smaller than $\deg G$, forcing
every derivative to be zero. In characteristic zero this would
make $G$ constant, a contradiction. Thus $e\ge2$, and
$\deg P\ge2\deg G\ge8$.
\end{proof}

\begin{corollary}\label{lb:quartic-sign}
A three-coordinate quartic defining an infinite subset of
$\OKS$ must take both positive and negative values at every
real embedding of $K$. In particular, a quartic over $\mathbb Q$
whose rational input projection is $\mathbb Z$ cannot be
nonnegative on $\mathbb R^3$.
\end{corollary}
\begin{proof}
Otherwise the real polynomial has a constant sign, and either
$P$ or $-P$ contradicts Theorem~\ref{lb:nonnegative-eight}.
\end{proof}

\begin{lemma}\label{lb:quartic-singular-surface}
Let $X\subset\PP^4_K$ be a geometrically integral quartic
hypersurface. Every geometric surface component of its singular
locus has degree at most three.
\end{lemma}
\begin{proof}
Work over an algebraic closure of $K$, and let $Y$ be such a
component. By Bertini's irreducibility theorem, applied twice,
a general plane $\Lambda\simeq\PP^2$ cuts $X$ in an integral
plane quartic $C$; see \cite[Lemma~37.32.3]{StacksBertini}.
The same general plane meets the reduced surface $Y$
transversely at $\deg Y$ distinct points. At each such point
all partial derivatives of the quartic equation of $X$ vanish,
so the restricted plane equation is singular there.

The arithmetic genus of a plane quartic is
$(4-1)(4-2)/2=3$. If $\widetilde C$ is its normalization, the
normalization exact sequence gives
\[
 3=g(\widetilde C)+\sum_{P\in\operatorname{Sing}(C)}\delta_P,
 \qquad
 \delta_P=\dim_{\overline K}
              (\widetilde{\mathcal O}_{C,P}/\mathcal O_{C,P}).
\]
Every singular point of an integral plane curve has
$\delta_P\ge1$, and the geometric genus is nonnegative.
Thus $C$ has at most three singular points, and $\deg Y\le3$.
\end{proof}

\begin{theorem}\label{lb:nonnegative-quartic-five}
Let $K$ have a real embedding $\sigma$ and let
$P\in K[T,Y_1,Y_2,Y_3]$ have degree at most four. Suppose that
$P^\sigma$ is nonnegative on $\mathbb R^4$. If the rational
input projection of $P$ is contained in $\OKS$, it is finite.
Consequently a nonnegative quartic with infinite $S$-integral
input projection needs at least five total coordinates.
\end{theorem}
\begin{proof}
The zero polynomial is excluded by the assumed projection
containment, while a nonzero constant has empty projection.
If all irreducible factors of $P$ over $K$ have degree at most
three, each factor has its rational zero set contained in that
of $P$, so Theorem~\ref{lb:cubic} gives finite input image for
each. It therefore suffices to consider an irreducible quartic.

Suppose first that the quartic is not absolutely irreducible.
In characteristic zero its geometric factors are distinct,
have equal degrees, and form one Galois orbit. Thus it is a
nonzero constant times either the norm of a quadratic polynomial
over a quadratic extension or the norm of a linear polynomial
over a degree-four extension. To see the precise field of
definition, normalize one nonzero coefficient of a geometric
factor to one; its stabilizer then fixes all its coefficients,
and its orbit has respectively two or four elements. On a
$K$-rational tuple the norm vanishes exactly when that polynomial
vanishes. Expanding in a $K$-basis of the extension gives
respectively two quadratic equations or four linear equations
over $K$. Their input image is finite by
Theorem~\ref{lb:quadratic-five} or
Corollary~\ref{lb:linear-systems}.

Finally suppose that $P$ is absolutely irreducible, and let
$X\subset\PP^4_K$ be the quartic obtained by homogenization.
Every rational affine zero of $P$ is a real minimum after
applying $\sigma$. All its first derivatives vanish. Euler's
identity for the homogenized polynomial shows that its derivative
in the homogenizing coordinate also vanishes. Therefore all
rational affine zeros belong to $\operatorname{Sing}(X)$.

This singular locus has dimension at most two. Components of
dimension at most one contribute finite input image by
Proposition~\ref{lb:curves}. If a surface component is not
geometrically integral, its rational points lie in a closed
subset of dimension at most one, with the same conclusion.
Every geometrically integral surface component has degree at
most three by Lemma~\ref{lb:quartic-singular-surface}; its
input image on the affine chart is finite by
Lemma~\ref{lb:small-degree-surface}. Taking the finite union
over components proves the theorem.
\end{proof}

\section{Scope of the coordinate bounds}

\begin{corollary}\label{lb:degree-sharpness}
Let $A\subseteq\OKS$ be infinite. If $A$ has a single-polynomial
presentation of degree four over $K$, then four is the least
possible degree of any such presentation, regardless of how many
witness coordinates are allowed. If $A$ has a presentation by a
system of quadratic equations, then two is the least possible
maximum degree of the equations in a system.
\end{corollary}
\begin{proof}
The first assertion follows from Theorem~\ref{lb:main}; the second
follows from Corollary~\ref{lb:linear-systems}.
\end{proof}

The bounds for systems can be displayed as follows; the degree
in the table is the largest degree among the equations.
\[
\begin{array}{c|c}
\text{maximum equation degree}&\text{necessary total coordinates}\\
\hline
1&\text{no finite number suffices}\\
2&\text{at least }5\\
3&\text{at least }4\\
4\text{ or greater}&\text{at least }3
\end{array}
\]
The nonzero conditions may have arbitrary degrees. For a single
equation, degrees at most three are impossible regardless of the
number of coordinates. All nonnegative quartics, including
positive sums of squares of quadratics, need at least five
coordinates, while arbitrary quartics have
only the unconditional three-coordinate lower bound proved here.
Nonnegative polynomials in three coordinates need degree at
least eight.

Let $\nu$ be the least total number of coordinates in a
single-polynomial definition of $\mathbb Z$ over $\mathbb Q$,
if such a definition exists, and put $\nu=+\infty$ otherwise.
The unconditional conclusion is
\[
                              \nu\ge3.
\]
An upper bound obtained from a construction is subject to that
construction's stated arithmetic hypotheses. The proofs above
do not establish $\nu=3$, do not exclude $\nu=3$, and do not
provide an attained minimum for quartics or quadratic systems.
In particular, the sharp lower bound on \emph{degree} supplied
by a valid quartic construction does not make its number of
coordinates optimal.

The distinction between the number of quantified coordinates
and the dimension of the witness fibers is also essential.
For terminology and general results concerning existential rank
and essential fiber dimension, see \cite{DDF}. In the convention
of the present paper, the total coordinate count is one more
than the number of existentially quantified coordinates.

\begin{thebibliography}{9}
\bibitem{Kollar}
J.~Koll\'ar, Unirationality of cubic hypersurfaces,
\emph{Journal of the Institute of Mathematics of Jussieu}
\textbf{1} (2002), 467--476, Theorem~2.
\url{https://arxiv.org/abs/math/0005146}.

\bibitem{Lang}
S.~Lang, \emph{Fundamentals of Diophantine Geometry},
Springer, 1983, Chapter~8, ``Siegel's Theorem and Integral Points'',
188--224.
\url{https://doi.org/10.1007/978-1-4757-1810-2_8}.

\bibitem{EGHP}
D.~Eisenbud, M.~Green, K.~Hulek, and S.~Popescu,
Small schemes and varieties of minimal degree,
\emph{American Journal of Mathematics} \textbf{128} (2006),
1363--1389; Introduction, inequality $(*)$.
\url{https://arxiv.org/abs/math/0404517}.

\bibitem{DDF}
N.~Daans, P.~Dittmann, and A.~Fehm,
Existential rank and essential dimension of diophantine sets,
arXiv:2102.06941v5, 2025.
\url{https://arxiv.org/abs/2102.06941}.

\bibitem{Mumford}
D.~Mumford, \emph{Algebraic Geometry I: Complex Projective Varieties},
Grundlehren der mathematischen Wissenschaften, vol.~221,
Springer, 1976, Theorem~5.11; corrected reprint, 1995.
\url{https://link.springer.com/book/9783540586579}.

\bibitem{StacksBertini}
The Stacks Project Authors, \emph{The Stacks Project},
Section~37.32, ``Bertini theorems'', Lemma~37.32.3.
\url{https://stacks.math.columbia.edu/tag/0G4C}.
\end{thebibliography}
\end{document}
